% Рисунок с доски: исходные модели (Кельвин — Фойгт, Максвелл) и две формы
% их объединения — стандартное линейное тело.
\begin{figure}[htbp]
\centering
\begin{tikzpicture}[>={Latex[length=4pt]}, scale=0.9, transform shape]
  % --- Кельвин — Фойгт ---
  \begin{scope}
    \node[font=\footnotesize, anchor=south] at (0.9,2.3) {материал К--Ф};
    \pic[rotate=90] at (0,0) {spring};
    \pic[rotate=90] at (1.8,0) {damper};
    \draw[thick] (0,0.9) -- (0,1.3) -- (1.8,1.3) -- (1.8,0.9);
    \draw[thick] (0,-0.9) -- (0,-1.3) -- (1.8,-1.3) -- (1.8,-0.9);
    \draw[thick] (0.9,1.3) -- (0.9,1.8);
    \draw[thick] (0.9,-1.3) -- (0.9,-1.8);
    \node[font=\footnotesize, anchor=north] at (0.9,-2.0)
      {$\Sdev = 2\mu\,\dev + 2\eta\,\dt{\dev}$};
  \end{scope}

  % --- Максвелл ---
  \begin{scope}[xshift=6cm]
    \node[font=\footnotesize, anchor=south] at (0,2.3) {Максвелл};
    \draw[thick] (0,1.8) -- (0,1.5);
    \pic[rotate=90] at (0,0.6) {spring};
    \draw[thick] (0,-0.3) -- (0,-0.3);
    \pic[rotate=90] at (0,-1.2) {damper};
    \draw[thick] (0,-2.1) -- (0,-2.4);
    \node[font=\footnotesize, anchor=north] at (0,-2.6)
      {$\dt{\dev} = \dfrac{\dt{\Sdev}}{2\mu} + \dfrac{\Sdev}{2\eta}$};
  \end{scope}
\end{tikzpicture}

\bigskip
\textit{Хотим их объединить.}

\bigskip
\begin{tikzpicture}[>={Latex[length=4pt]}, scale=0.9, transform shape]
  % --- форма 1: ветвь Максвелла параллельно пружине ---
  \begin{scope}
    \pic[rotate=90] at (0,1.0) {spring};
    \node[font=\footnotesize, anchor=east] at (-0.5,1.0) {$\mu$};
    \pic[rotate=90] at (0,-1.0) {damper};
    \node[font=\footnotesize, anchor=east] at (-0.5,-1.0) {$\eta$};
    \pic[rotate=90] at (2.2,0) {spring};
    \node[font=\footnotesize, anchor=west] at (2.7,0) {$\mu_1$};

    \draw[thick] (0,1.9) -- (0,2.3) -- (2.2,2.3) -- (2.2,0.9);
    \draw[thick] (0,-1.9) -- (0,-2.3) -- (2.2,-2.3) -- (2.2,-0.9);
    \draw[thick] (1.1,2.3) -- (1.1,2.8);
    \draw[thick] (1.1,-2.3) -- (1.1,-2.8);
  \end{scope}

  % --- форма 2: пружина последовательно со звеном К — Ф ---
  \begin{scope}[xshift=7.5cm]
    \draw[thick] (1.1,2.8) -- (1.1,2.6);
    \pic[rotate=90] at (1.1,1.7) {spring};
    \draw[thick] (1.1,0.8) -- (1.1,0.6);
    \pic[rotate=90] at (0,-0.4) {damper};
    \pic[rotate=90] at (2.2,-0.4) {spring};
    \draw[thick] (0,0.5) -- (0,0.6) -- (2.2,0.6) -- (2.2,0.5);
    \draw[thick] (0,-1.3) -- (0,-1.5) -- (2.2,-1.5) -- (2.2,-1.3);
    \draw[thick] (1.1,-1.5) -- (1.1,-2.8);
  \end{scope}
\end{tikzpicture}
\caption{Объединение моделей: ветвь Максвелла параллельно пружине $\mu_1$
(слева) и пружина последовательно со звеном Кельвина–Фойгта (справа). Обе схемы
дают стандартное линейное тело.}
\label{fig:standard-combine}
\end{figure}
